% ============================================================
%  Hardware--Software Co-design for Autonomous Vehicle Control
%  via Cycle-Accurate Simulation
%  IEEE CDC style — double-column, 10pt
% ============================================================
\documentclass[conference]{IEEEtran}

% ---- Packages -----------------------------------------------
\usepackage{amsmath,amssymb,amsthm}
\usepackage{bm}
\usepackage{graphicx}
\usepackage{booktabs}
\usepackage{multirow}
\usepackage{xcolor}
\usepackage{hyperref}
\usepackage{cite}
\usepackage[T1]{fontenc}

% ---- Macros --------------------------------------------------
\newcommand{\R}{\mathbb{R}}
\newcommand{\N}{\mathbb{N}}
\newcommand{\RMSE}{\mathrm{RMSE}}
\newcommand{\calX}{\mathcal{X}}
\newcommand{\calU}{\mathcal{U}}
\newcommand{\calS}{\mathcal{S}}
\newcommand{\calH}{\mathcal{H}}
\newcommand{\calG}{\mathcal{G}}
\newcommand{\TODO}[1]{\textcolor{red}{\textbf{[TODO: #1]}}}

% ---- Theorem environments -----------------------------------
\newtheorem{theorem}{Theorem}
\newtheorem{proposition}{Proposition}
\newtheorem{lemma}{Lemma}
\newtheorem{corollary}{Corollary}
\newtheorem{definition}{Definition}
\newtheorem{remark}{Remark}
\newtheorem{assumption}{Assumption}

% =============================================================
\begin{document}

\title{Hardware--Software Co-design for Autonomous Vehicle Control\\
  via Cycle-Accurate Simulation}

\author{\IEEEauthorblockN{Author~1\IEEEauthorrefmark{1},
    Author~2\IEEEauthorrefmark{2},
    Author~3\IEEEauthorrefmark{1}}
  \IEEEauthorblockA{\IEEEauthorrefmark{1}Affiliation~1 \quad
    \IEEEauthorrefmark{2}Affiliation~2\\
    \texttt{\{email1, email2\}@placeholder.edu}}
}

\maketitle

% =============================================================
\begin{abstract}
  In autonomous driving, control algorithms ranging from PID to
  nonlinear model predictive control (NMPC) must execute within a
  strict sample period on embedded processors whose computational
  resources vary widely.  The computation time depends on the
  microarchitecture---clock speed, cache size, pipeline depth---and
  interacts with the controller's algorithmic complexity in ways that
  are difficult to predict analytically.  This paper presents a
  co-simulation framework coupling \emph{SHARC} (Simulator for
  Hardware Architecture and Real-time Control) with the \emph{CARLA}
  driving simulator, enabling systematic evaluation of how hardware
  and controller design choices jointly affect closed-loop safety.  We
  formalize the delay-induced closed-loop dynamics, characterize
  conditions under which computational deadline misses cascade, and
  design an experimental protocol that sweeps hardware parameters
  (clock frequency, cache size) and software parameters (MPC
  prediction horizon, sample period) across safety-critical driving
  scenarios.
\end{abstract}

\begin{IEEEkeywords}
  Hardware--software co-design, autonomous driving, model predictive
  control, computational delay, microarchitecture simulation.
\end{IEEEkeywords}

% =============================================================
\section{Introduction}\label{sec:intro}

Consider an autonomous vehicle executing a model predictive controller
at 10\,Hz on an embedded processor.  At every sample instant, the
controller receives the vehicle state, solves an optimization problem,
and outputs throttle, steering, and brake commands.  The time $\tau$
this computation takes is \emph{not} a constant---it depends on the
optimization landscape (which changes with the vehicle state), the
processor's clock frequency, its cache hierarchy, and many
microarchitectural details.  When $\tau$ exceeds the sample
period~$T$, the controller misses its deadline: the plant continues to
apply a stale control input.  In safety-critical scenarios---emergency
braking, intersection crossing---a single missed deadline can lead to
collision.

Despite this tight coupling, control algorithm design and hardware
selection are typically treated independently.  Control engineers assume
instantaneous computation; chip designers optimize for generic
workloads.  When timing violations surface during integration testing,
the only recourses are expensive hardware upgrades or ad-hoc algorithm
simplification.

\emph{What is needed is a simulation tool that evaluates the
closed-loop system---physics, control law, and computational
hardware---simultaneously.}  This paper provides such a tool by
integrating two existing simulators:

\begin{itemize}
  \item \textbf{SHARC}~\cite{wintz2025sharc}: uses the Scarab
    cycle-accurate microarchitecture simulator~\cite{scarab} to
    measure the execution time of a controller binary on a specified
    (possibly hypothetical) processor.
  \item \textbf{CARLA}~\cite{dosovitskiy2017carla}: an open-source,
    high-fidelity driving simulator with realistic vehicle dynamics,
    road networks, sensors, and traffic.
\end{itemize}

\noindent\textbf{Contributions.}
\begin{enumerate}
  \item A \textbf{CARLA dynamics interface} for SHARC that replaces
    abstract ODE integrators with high-fidelity, synchronous vehicle
    simulation (Section~\ref{sec:carla}).
  \item \textbf{Formal analysis} of delay-induced closed-loop
    degradation: we define the \emph{hardware--performance map} and
    state conditions under which a single deadline miss triggers
    cascading error growth (Section~\ref{sec:analysis}).
  \item An \textbf{experimental co-design protocol} that jointly sweeps
    hardware parameters (clock frequency, cache size) and software
    parameters (MPC horizon $N_p$, sample period $T$) across
    safety-critical driving scenarios
    (Section~\ref{sec:experiments}).
\end{enumerate}

% =============================================================
\section{Modeling Framework}\label{sec:model}

We review the SHARC modeling framework~\cite{wintz2025sharc} and
introduce the notation used throughout.

\subsection{Plant Dynamics}

A plant with state $x \in \R^{n_x}$, control $u \in \R^{n_u}$,
exogenous input $w \in \R^{n_w}$, and output $y \in \R^{n_y}$ evolves
in continuous time as
\begin{equation}\label{eq:ct}
  \dot{x} = \tilde{f}(t, x, u, w), \qquad y = h(x, u, w).
\end{equation}
SHARC discretizes~\eqref{eq:ct} on the grid $t_k := kT$, $k \in \N$,
with sample period $T > 0$:
\begin{equation}\label{eq:disc}
  x_{k+1} = f(t_k, x_k, u_k, w_k),
\end{equation}
where $f$ is obtained by numerical integration under zero-order hold
(ZOH) on $u$, and $w_k := w(t_k)$.  Users may alternatively supply $f$
directly (for hybrid systems, stochastic models, or---as in our
case---external simulators).

\subsection{Computation Delay}\label{sec:delay_model}

A controller evaluates
$\tilde{u}_{k+1} := g(t_k, y_k) \in \R^{n_u}$
at each sample time $t_k$.  The computation requires $\tau_k > 0$
seconds, so the result is available at time
$\hat{t}_{k+1} := t_k + \tau_k$.

\begin{definition}[Delay-Aware ZOH]\label{def:zoh}
  The applied input under computation delay is
  \begin{equation}\label{eq:zoh_delay}
    u(t) = \begin{cases}
      u_k & t \in [t_k,\, t_{k+1}),
        \;\;\text{if } \tau_k \geq T, \\[2pt]
      u_k & t \in [t_k,\, \hat{t}_{k+1}),
        \;\;\text{if } \tau_k < T, \\[2pt]
      \tilde{u}_{k+1} & t \in [\hat{t}_{k+1},\, t_{k+1}),
        \;\;\text{if } \tau_k < T.
    \end{cases}
  \end{equation}
\end{definition}

\begin{definition}[Deadline Miss]\label{def:miss}
  A \emph{deadline miss} at step~$k$ occurs when $\tau_k \geq T$.
  The \emph{miss rate} over $K$ steps is
  $\rho := \frac{1}{K}\sum_{k=0}^{K-1}\mathbf{1}[\tau_k \geq T]$.
\end{definition}

\begin{definition}[Effective Closed-Loop Dynamics]\label{def:cl}
  Let $d_k := \min(\tau_k, T)$.
  The closed-loop system under delayed ZOH is
  \begin{equation}\label{eq:cl}
    x_{k+1} = f\!\bigl(t_k,\, x_k,\, \phi(u_k, \tilde{u}_{k+1},
    d_k, T),\, w_k\bigr),
  \end{equation}
  where $\phi$ is the delay-interpolated input:
  \begin{equation}\label{eq:phi}
    \phi(u, \tilde{u}, d, T) :=
    \frac{d}{T}\,u + \frac{T - d}{T}\,\tilde{u}.
  \end{equation}
\end{definition}

\begin{remark}
  When $\tau_k \geq T$, $d_k = T$ and $\phi = u_k$: the full period
  uses the stale input.  When $\tau_k = 0$, $\phi = \tilde{u}_{k+1}$:
  ideal operation.  Expression~\eqref{eq:phi} continuously
  interpolates these extremes.
\end{remark}

\subsection{Microarchitecture-Dependent Computation Time}

The computation delay is
\begin{equation}\label{eq:tau_model}
  \tau_k = \mathcal{T}(g,\, x_k,\, \calH),
\end{equation}
where $g$ is the controller, $x_k$ the current state, and
$\calH := (f_\text{clk}, S_\text{L1}, S_\text{L2}, \ldots)
\in \R_+^{n_h}$ is a \emph{hardware configuration vector} specifying
clock frequency, cache sizes, pipeline width, etc.  In SHARC,
$\mathcal{T}$ is not modeled analytically---it is \emph{evaluated} by
the Scarab microarchitecture simulator~\cite{scarab}, which executes
the controller binary instruction-by-instruction on the target hardware
model and returns a cycle-accurate execution time.  Hardware is
specified via configuration files: changing one parameter (say, halving
the L1 cache) is a single-line edit, making large sweeps practical.

\begin{definition}[Hardware \& Software Configuration Spaces]\label{def:spaces}
  The \emph{hardware space}
  $\calH \subset \R_+^{n_h}$ and the \emph{software space}
  $\calG := \{\theta = (N_p, N_c, T, Q, R, \ldots)\}
  \subset \R^{n_s}$ parameterize, respectively, the processor
  microarchitecture and the controller algorithm.
\end{definition}

% =============================================================
\section{CARLA as Plant Dynamics}\label{sec:carla}

SHARC provides a modular dynamics interface: any class that implements
$f(t_k, x_k, u_k, w_k) \to x_{k+1}$ can serve as the plant.
CARLA~\cite{dosovitskiy2017carla} is a natural fit because its
synchronous mode lets an external client control \emph{exactly when}
time advances---precisely the requirement of~\eqref{eq:disc}.

Thanks to SHARC's user-friendly, extensible design, integrating CARLA
requires implementing only a single dynamics class that (i)~connects to
the CARLA server and sets the physics time-step to $T$, (ii)~applies
the control input $u_k$ to the ego vehicle, (iii)~ticks the world
forward, and (iv)~reads back the new state $x_{k+1}$.  No modification
to SHARC's core codebase is needed.

\subsection{State and Control Representation}

For the ego vehicle:
\begin{align}
  x_k &:= (p_x,\, p_y,\, \psi,\, v,\, w_x^+,\,
    w_y^+)^\top \in \R^6, \label{eq:state}\\
  u_k &:= (\alpha,\, \delta,\, \beta)^\top \in
    [0, 1] \times [-1, 1] \times [0, 1], \label{eq:control}
\end{align}
where $(p_x, p_y)$ is position, $\psi$ yaw, $v$ speed,
$(w_x^+, w_y^+)$ the next waypoint from the CARLA HD map,
and $(\alpha, \delta, \beta)$ are throttle, steering, brake.

The transition $x_{k+1} = F_{\mathrm{CARLA}}(x_k, u_k, w_k)$ is
evaluated by CARLA's internal physics engine, which models tire
forces, aerodynamic drag, engine torque curves, suspension, and road
geometry---providing dynamics fidelity far beyond closed-form ODEs.
From the SHARC framework's perspective, $F_{\mathrm{CARLA}}$ is simply
a concrete instantiation of $f$ in~\eqref{eq:disc}.

% =============================================================
\section{Controller Formulations}\label{sec:controllers}

We study three controllers spanning the computational complexity
spectrum.  Each is compiled as a C++ binary whose execution is
instrumented by Scarab to measure $\tau_k$.

\subsection{PID Controller}\label{sec:pid}

Vehicle control is decomposed into longitudinal and lateral PID loops.

\noindent\textbf{Longitudinal (speed tracking).}
\begin{equation}\label{eq:pid_lon}
  e_k^v := v^* - v_k, \quad
  \sigma_k = K_P^\ell\, e_k^v +
  K_I^\ell \!\sum_{j=0}^{k} e_j^v T +
  K_D^\ell\, \frac{e_k^v - e_{k-1}^v}{T}.
\end{equation}
Throttle and brake are set as
$\alpha_k = \mathrm{clip}(\sigma_k,\, 0,\, \alpha_\text{max})$,
$\beta_k = \mathrm{clip}(-\sigma_k,\, 0,\, \beta_\text{max})$.

\noindent\textbf{Lateral (waypoint tracking).}
With heading error
$e_k^\psi := \mathrm{atan2}(w_y^+ - p_y,\, w_x^+ - p_x) - \psi_k$,
\begin{equation}\label{eq:pid_lat}
  \delta_k = \mathrm{clip}\!\Bigl(K_P^r e_k^\psi + K_I^r
  \textstyle\sum_{j} e_j^\psi T + K_D^r
  \frac{e_k^\psi - e_{k-1}^\psi}{T},\;
  \pm\delta_\text{max}\Bigr).
\end{equation}

PID requires $O(1)$ operations per step; we expect
$\tau_k^{\mathrm{PID}} \ll T$ on virtually all hardware.

\subsection{Linear MPC (LMPC)}\label{sec:lmpc}

We linearize the vehicle dynamics around a reference
$(\bar{x}_k, \bar{u}_k)$:
\begin{equation}\label{eq:lin}
  \Delta x_{k+1} = A_k\,\Delta x_k + B_k\,\Delta u_k, \quad
  A_k := \tfrac{\partial f}{\partial x}\Big|_{\bar{x}_k},\;
  B_k := \tfrac{\partial f}{\partial u}\Big|_{\bar{u}_k}.
\end{equation}
At each $t_{k_0}$, LMPC solves the convex QP:
\begin{subequations}\label{eq:lmpc}
\begin{align}
  \min_{\Delta u_{k_0:k_0+N_c-1}} \;\;&
    \sum_{k=k_0}^{k_0+N_p-1}\bigl(
    \|\Delta x_{k|k_0}\|_Q^2 + \|\Delta u_{k|k_0}\|_R^2\bigr)\\
  \text{s.t.}\;\; & \Delta x_{k+1|k_0} = A_k\Delta x_{k|k_0}
    + B_k\Delta u_{k|k_0}, \\
  & u_{k|k_0} \in \calU,\quad x_{k|k_0} \in \calX,
\end{align}
\end{subequations}
with prediction horizon $N_p$, control horizon $N_c \leq N_p$,
$Q \succeq 0$, $R \succ 0$.

\subsection{Nonlinear MPC (NLMPC)}\label{sec:nlmpc}

NLMPC uses a kinematic bicycle model $F_\text{bike}$ for prediction:
\begin{subequations}\label{eq:nlmpc}
\begin{align}
  \min_{u_{k_0:k_0+N_c-1}} \;\;&
    \sum_{k=k_0}^{k_0+N_p-1}
    \|x_{k|k_0} - x_k^\text{ref}\|_Q^2 + \|u_{k|k_0}\|_R^2\\
  \text{s.t.}\;\;& x_{k+1|k_0} = F_\text{bike}(x_{k|k_0},
    u_{k|k_0}), \\
  & u_{k|k_0} \in \calU,\quad x_{k|k_0} \in \calX.
\end{align}
\end{subequations}
This nonconvex NLP is solved iteratively via
SLSQP~\cite{kraft1988slsqp}.

\begin{remark}[Complexity ordering]\label{rem:complexity}
  $\tau_k^{\mathrm{PID}} \ll \tau_k^{\mathrm{LMPC}} \ll
  \tau_k^{\mathrm{NLMPC}}$:
  PID is $O(1)$; LMPC solves a QP of dimension
  $O(N_c n_u)$; NLMPC iteratively solves a nonconvex NLP.
  This ordering makes the three controllers ideal probes for the
  effect of hardware constraints on closed-loop performance.
\end{remark}

% =============================================================
\section{Delay-Induced Degradation Analysis}\label{sec:analysis}

We now formalize the mechanisms by which computational delays degrade
closed-loop performance.

\subsection{The Hardware--Performance Map}

\begin{definition}[Hardware--Performance Map]\label{def:hp_map}
  Fix a controller $g_\theta$ with software parameters $\theta$,
  a scenario $\calS$ (initial conditions, reference, disturbances),
  and a performance metric $J : \{x_k\}_{k=0}^K \to \R_{\geq 0}$.
  The \emph{hardware--performance map} is
  \begin{equation}\label{eq:hp_map}
    \Phi_{\theta, \calS} : \calH \to \R_{\geq 0},
    \quad h \mapsto J\!\bigl(\{x_k(h)\}_{k=0}^K\bigr),
  \end{equation}
  where $\{x_k(h)\}$ is the trajectory generated under hardware $h$
  via~\eqref{eq:cl}.
\end{definition}

The map $\Phi$ encodes the central co-design question: \emph{for a
given controller and driving scenario, how does closed-loop cost change
as we modify the processor?}

\begin{definition}[Minimum Viable Hardware]\label{def:mvh}
  Given a performance threshold $J^* > 0$, the
  \emph{minimum viable hardware set} is
  \begin{equation}\label{eq:mvh}
    \calH^*(\theta, \calS, J^*) :=
    \bigl\{h \in \calH : \Phi_{\theta, \calS}(h) \leq J^*\bigr\}.
  \end{equation}
  When $\calH$ is parameterized by clock frequency alone,
  $f_{\min}^* := \inf\{f_\text{clk}: (f_\text{clk}, \cdot) \in
  \calH^*\}$ is the \emph{minimum viable clock frequency}.
\end{definition}

\subsection{Cascading Deadline Misses}

The most dangerous failure mode occurs when misses self-reinforce.

\begin{assumption}\label{asm:tau_lipschitz}
  The computation time is locally Lipschitz in the state: there exists
  $L_\tau > 0$ such that
  $|\tau_k - \tau_j| \leq L_\tau \|x_k - x_j\|$ for all
  $x_k, x_j$ in a neighborhood $\calX_0$ of the reference trajectory.
\end{assumption}

This is natural for iterative solvers: larger state deviations require
more solver iterations, increasing execution time.

\begin{proposition}[Cascading Condition]\label{prop:cascade}
  Under Assumption~\ref{asm:tau_lipschitz}, let $L_f$ be the Lipschitz
  constant of $f$ with respect to its input argument, and let
  $e_{k} := \|x_{k} - \bar{x}_{k}\|$ be the tracking error.
  If a deadline miss occurs at step~$k$ ($\tau_k \geq T$, so
  $\phi = u_k$) and
  \begin{equation}\label{eq:cascade_cond}
    L_\tau \cdot L_f \cdot \|u_k - \tilde{u}_{k+1}\| > 1,
  \end{equation}
  then the miss at step $k$ increases the tracking error, which
  in turn increases $\tau_{k+1}$, creating a positive feedback loop.
\end{proposition}

\begin{proof}
  A deadline miss forces $\phi = u_k$.  The one-step error grows as
  \begin{equation*}
    e_{k+1} \leq e_k + L_f\,\|u_k - \tilde{u}_{k+1}\|.
  \end{equation*}
  By Assumption~\ref{asm:tau_lipschitz},
  $\tau_{k+1} \geq \tau_k + L_\tau(e_{k+1} - e_k)
  \geq \tau_k + L_\tau L_f \|u_k - \tilde{u}_{k+1}\|$.
  When~\eqref{eq:cascade_cond} holds,
  $\tau_{k+1} > \tau_k \geq T$, so step $k{+}1$ is also a miss,
  and $e_{k+2}$ grows further.
\end{proof}

\begin{corollary}[PID Immunity]\label{cor:pid}
  For PID, $\tau_k^{\mathrm{PID}}$ is essentially constant
  ($L_\tau \approx 0$) since the computation is a fixed number of
  arithmetic operations independent of $x_k$.
  Thus~\eqref{eq:cascade_cond} is never satisfied, and PID does not
  exhibit cascading misses.
\end{corollary}

\begin{remark}[Implications for controller selection]
  Proposition~\ref{prop:cascade} reveals a fundamental tension:
  MPC yields better tracking when $\tau_k \ll T$ but is \emph{fragile}
  near the deadline boundary.  PID is suboptimal but \emph{robust}
  to hardware constraints.  This motivates co-design: choose the most
  complex controller the hardware can guarantee to run within~$T$.
\end{remark}

\subsection{The Co-design Optimization Problem}

\begin{definition}[Co-design Problem]\label{def:codesign}
  Given a scenario $\calS$, controller family $\{g_\theta\}$, and
  hardware cost $C_\text{hw}$ (chip area, power), the
  \emph{co-design problem} is
  \begin{subequations}\label{eq:codesign}
  \begin{align}
    \min_{\theta \in \calG,\; h \in \calH}\;\;&
      \Phi_{\theta,\calS}(h) \\
    \text{s.t.}\;\;& \rho(\theta, h, \calS) \leq \rho_\text{max},\\
    & C_\text{hw}(h) \leq C_\text{budget}.
  \end{align}
  \end{subequations}
\end{definition}

Problem~\eqref{eq:codesign} is intractable analytically: $\Phi$ is a
black-box evaluated through cycle-accurate simulation.  The
SHARC--CARLA framework provides the oracle for $\Phi$, enabling
systematic search over $(\theta, h)$.

% =============================================================
\section{Experimental Design}\label{sec:experiments}

We evaluate $\Phi_{\theta,\calS}(h)$ across a representative grid.

\subsection{Driving Scenarios ($\calS$)}

Table~\ref{tab:scenarios} lists six CARLA scenarios ordered by safety
criticality.  Scenarios with smaller $T$ demand faster computation.

\begin{table}[ht]
  \centering
  \caption{Driving scenarios.}
  \label{tab:scenarios}
  \footnotesize
  \renewcommand{\arraystretch}{1.15}
  \begin{tabular}{lp{2.5cm}c}
    \toprule
    \textbf{Scenario} & \textbf{Key challenge}
      & $T$ (s) \\
    \midrule
    Highway cruising       & Speed tracking        & 0.25 \\
    Lane following         & Lateral precision      & 0.25 \\
    Adaptive cruise ctrl   & Lead-car gap           & 0.20 \\
    Intersection crossing  & Collision avoidance    & 0.15 \\
    Emergency braking      & Min.\ stopping dist.   & 0.10 \\
    High-speed racing      & Stability at limit     & 0.10 \\
    \bottomrule
  \end{tabular}
\end{table}

\subsection{Hardware Sweep ($\calH$)}

We sweep two hardware dimensions:
\begin{enumerate}
  \item \textbf{Clock frequency}
    $f_\text{clk} \in \{250, 500, 625, 750, 1000, 1500, 2000\}$\,MHz.
  \item \textbf{L1 data cache}
    $S_\text{L1} \in \{8, 16, 32, 64, 128\}$\,kB.
\end{enumerate}
All other Scarab parameters are held at a baseline ARM Cortex-A
profile.  This yields a $7 \times 5 = 35$-point grid.

\subsection{Software Sweep ($\calG$)}

For MPC controllers, we sweep:
\begin{enumerate}
  \item $N_p \in \{5, 10, 15, 20, 30\}$ (prediction horizon).
  \item $N_c \in \{\lceil N_p/2\rceil, N_p\}$ (control horizon).
\end{enumerate}
PID gains are fixed; PID serves as a hardware-insensitive baseline.

\subsection{Performance Metrics}

\begin{definition}[Tracking Error]\label{def:rmse}
  $J_\text{pos} := \RMSE_\text{pos} =
  \sqrt{\frac{1}{K}\sum_{k=0}^{K-1}
  [(p_x^k - \bar{p}_x^k)^2 + (p_y^k - \bar{p}_y^k)^2]}$.
\end{definition}

\begin{definition}[Safety Violation Count]\label{def:safety}
  $J_\text{safe} := n_\text{col} + n_\text{lane} + n_\text{speed}$,
  counting collisions, lane departures ($|e_y| > 1.5$\,m), and
  speed exceedances ($v > 1.2\,v^*$).
\end{definition}

\noindent Additional metrics: $\rho$ (Def.~\ref{def:miss}),
$\bar{\tau} := \frac{1}{K}\sum_k \tau_k$, and delay distribution
$p(\tau_k)$.

\subsection{Planned Analyses}

\noindent\textbf{1.\ Delay distributions.}
Box plots of $\tau_k$ per controller$\times$hardware, annotated with
deadline $T$.

\noindent\textbf{2.\ Hardware--performance cliff.}
$J_\text{pos}$ vs.\ $f_\text{clk}$, identifying $f_\text{min}^*$
(Def.~\ref{def:mvh}) for each controller.

\noindent\textbf{3.\ Cascading miss dynamics.}
Time series of $(\tau_k, e_k)$ illustrating
Proposition~\ref{prop:cascade}.

\noindent\textbf{4.\ Co-design Pareto frontier.}
$J_\text{pos}$ vs.\ $f_\text{clk}$ for multiple $N_p$, identifying
Pareto-optimal $(N_p, f_\text{clk})$ pairs (Problem~\eqref{eq:codesign}).

\noindent\textbf{5.\ Cache sensitivity.}
$J_\text{pos}$ vs.\ $S_\text{L1}$ at fixed clock, revealing whether
cache is a first- or second-order effect.

\noindent\textbf{6.\ Safety heatmap.}
$J_\text{safe}$ on the $(f_\text{clk}, N_p)$ grid for emergency
braking and intersection scenarios.

% =============================================================
\section{Results}\label{sec:results}

\TODO{All experimental results are pending.  The following subsections
  will be populated once experiments are complete.}

\subsection{Delay Distributions}
\TODO{Box plots and delay tables.}

\subsection{Hardware--Performance Cliff}
\TODO{$J_\text{pos}$ vs.\ $f_\text{clk}$ plots with cliff annotation.}

\subsection{Cascading Miss Dynamics}
\TODO{Time-series of $(\tau_k, e_k)$ per
  Proposition~\ref{prop:cascade}.}

\subsection{Co-design Pareto Analysis}
\TODO{Pareto frontier over $(N_p, f_\text{clk})$.}

\subsection{Cache Sensitivity}
\TODO{$J_\text{pos}$ vs.\ $S_\text{L1}$ curves.}

\subsection{Safety Event Analysis}
\TODO{Heatmaps of $J_\text{safe}$ on $(f_\text{clk}, N_p)$ grid.}

% =============================================================
\section{Discussion and Conclusion}\label{sec:conclusion}

\TODO{To be written after experimental results.  Key themes:
  (1)~hardware--performance cliff characterization,
  (2)~PID as robust fallback (Corollary~\ref{cor:pid}),
  (3)~MPC horizon as co-design lever,
  (4)~cache vs.\ clock importance,
  (5)~generalizability to other CPS,
  (6)~limitations (CARLA fidelity, single-agent),
  (7)~future work: multi-agent, control barrier functions,
  Bayesian optimization over $(\theta, h)$.}

% =============================================================
\bibliographystyle{IEEEtran}
\bibliography{references}

\end{document}
